\section{Pattern Validation: Classification Models}
\label{sec:3_8}

Classification models validate whether discovered patterns carry predictive signal, serving as the final stage of the pattern discovery pipeline.

\subsection{Baseline Models}

Two baseline models establish performance floors:

\textbf{Logistic Regression:} Linear baseline providing interpretable coefficients. Configuration uses L2 regularization (C=1.0) with class weight balancing. Logistic regression tests whether linear combinations of features suffice for fraud detection.

\textbf{Random Forest:} Tree ensemble baseline capturing non-linear interactions without graph information. Configuration:
\begin{itemize}
    \item Trees: 100
    \item Maximum depth: 10
    \item Class weight: balanced
\end{itemize}

Random Forest establishes the performance achievable with tabular features and non-linear modeling before graph augmentation.

\subsection{XGBoost Configuration}

XGBoost serves as the primary classification model due to demonstrated superiority on tabular data.

\textbf{Hyperparameter Optimization:} Optuna TPE (Tree-structured Parzen Estimator) sampler with 50 trials optimizes for AUC-PR:

\begin{table}[h]
\centering
\begin{tabular}{lll}
\hline
Parameter & Search Range & Best Value \\
\hline
\texttt{n\_estimators} & [50, 500] & 250 \\
\texttt{max\_depth} & [3, 12] & 6 \\
\texttt{learning\_rate} & [0.01, 0.3] & 0.1 \\
\texttt{min\_child\_weight} & [1, 10] & 3 \\
\texttt{subsample} & [0.6, 1.0] & 0.8 \\
\texttt{colsample\_bytree} & [0.6, 1.0] & 0.8 \\
\texttt{scale\_pos\_weight} & [1.0, 20.0] & 11.0 \\
\hline
\end{tabular}
\caption{XGBoost hyperparameter optimization results.}
\end{table}

\textbf{Class Imbalance Handling:} Rather than synthetic oversampling (SMOTE), we use \texttt{scale\_pos\_weight} to adjust loss weighting. Setting \texttt{scale\_pos\_weight} $\approx$ 11 (matching the 1:11 imbalance ratio) increases the penalty for false negatives on fraudulent listings. This approach avoids synthetic sample artifacts and integrates naturally with gradient boosting.

\subsection{Hybrid Pipeline}

The hybrid GNN+XGBoost pipeline operates in two stages:

\textbf{Stage 1 -- Pattern Extraction:}
\begin{itemize}
    \item Train GNN on heterogeneous graph.
    \item Generate 16-dimensional embeddings for all listings.
\end{itemize}

\textbf{Stage 2 -- Pattern Validation:}
\begin{itemize}
    \item Concatenate embeddings with tabular features: $\mathbf{f}_v = [\mathbf{x}_v \| \mathbf{z}_v]$
    \item Train XGBoost on enriched 81-dimensional feature vectors.
    \item Evaluate whether relational pattern representations improve classification.
\end{itemize}

This separation enables:
\begin{itemize}
    \item Independent debugging of GNN and classifier components.
    \item Embedding caching for efficient inference.
    \item Direct SHAP analysis on the final classifier.
    \item Quantification of GNN embedding contribution to predictions.
\end{itemize}
